% DocShield AI — Talent References
% SUM 72H Build Challenge — Deliverable 05
% Compile: pdflatex 05-talent-references.tex  (twice)

\documentclass[10pt,a4paper]{article}

\usepackage[T1]{fontenc}
\usepackage[utf8]{inputenc}
\usepackage[margin=16mm,top=14mm,bottom=12mm]{geometry}

\usepackage{charter}
\usepackage[activate={true,nocompatibility},final,tracking=true,
            factor=1100,stretch=12,shrink=12]{microtype}

\usepackage{amssymb}
\usepackage{xcolor}
\usepackage{enumitem}
\usepackage[most]{tcolorbox}
\usepackage[hidelinks]{hyperref}

\definecolor{navy}{HTML}{123A5F}
\definecolor{steel}{HTML}{3E6D9C}
\definecolor{good}{HTML}{1B6E4F}
\definecolor{goodbg}{HTML}{F1F8F4}
\definecolor{ink}{HTML}{16202A}
\definecolor{muted}{HTML}{68727E}

\pagestyle{empty}
\setlength{\parindent}{0pt}
\setlength{\parskip}{3pt}
\color{ink}

%% ------------------------------------------------------ encadrés à onglet
\tcbset{
  callout/.style={
    enhanced, sharp corners=downhill, arc=1.5pt,
    colback=white, boxrule=0.8pt,
    left=6pt, right=6pt, top=8pt, bottom=6pt,
    attach boxed title to top left={xshift=7pt, yshift=-\tcboxedtitleheight/2},
    boxed title style={arc=1pt, boxrule=0pt, left=5pt, right=5pt,
                       top=1.5pt, bottom=1.5pt},
    fonttitle=\footnotesize\bfseries\color{white},
  }
}

\newtcolorbox{note}[1]{callout, colframe=navy, colbacktitle=navy, title={#1}}
\newtcolorbox{person}[1]{callout, colframe=navy, colbacktitle=navy, title={#1}}
\newtcolorbox{gapbox}[1]{callout, colframe=good, colback=goodbg,
  colbacktitle=good, title={#1}}

\newcommand{\heading}[2]{%
  \par\vspace{8pt}{\color{navy}\bfseries #1\hspace{6pt}#2}\par\vspace{-3pt}}

% Libellé en attaque de paragraphe : compact, et sans le décalage vertical
% qu'une grille à deux colonnes introduit entre le libellé et sa valeur.
\newcommand{\fld}[1]{\textbf{\color{navy}#1}\hspace{5pt}}

\begin{document}

%% ------------------------------------------------------ en-tête
\noindent
\begin{minipage}[b]{0.62\linewidth}
  \colorbox{navy}{\makebox[7mm][c]{\rule[-2pt]{0pt}{12pt}\color{white}\bfseries 05}}%
  \hspace{5pt}{\Large\bfseries\color{navy}Talent References}\\[3pt]
  \hspace{12mm}{\color{muted}Who I would build DocShield AI with}
\end{minipage}\hfill
\begin{minipage}[b]{0.36\linewidth}
  \raggedleft\footnotesize\color{muted}
  SUM 72H Build Challenge\\[1pt]
  Deliverable 05 --- Talent References\\[1pt]
  Ahmed Hajji \ $\cdot$ \ Individual
\end{minipage}

\vspace{5pt}
\textcolor{navy}{\rule{\linewidth}{1pt}}

%% ------------------------------------------------------ cadrage
\heading{}{Selection logic}

I did not pick the strongest engineers I know. I picked the three people who
close the three weaknesses the 72-hour build actually exposed.

\vspace{2pt}
\begin{gapbox}{Gaps the prototype revealed}
\fld{Usability} Demo results were once mistaken for a real analysis because the
warning was not visible enough. A screening tool that misleads its own user fails.

\fld{Fraud expertise} An AI-generated invoice scored low risk. Knowing which
fraud patterns matter, and what a false positive costs, is domain knowledge I do
not have.

\fld{Production} No rate limiting, a payload ceiling on the host, no
persistence. The prototype runs; it is not yet a service.
\end{gapbox}

%% ------------------------------------------------------ 1
\heading{1.}{Product and user experience}

\begin{person}{Mohamed Sadik Derbi \ $\cdot$ \ Product / UX Engineer}
\fld{Contact} +212 651 843 140 \ $\cdot$ \
\href{https://www.linkedin.com/in/mohamed-sadikderbi/}{linkedin.com/in/mohamed-sadikderbi}

\fld{Key capability} Product thinking, UX/UI design, full-stack development.

\fld{Why him} DocShield produces a correct analysis that a reviewer can still
misread. Turning a risk score into a decision someone trusts is a product
problem, not a model problem --- and it is the part I am weakest on.

\fld{Role} Design the verification workflow, the evidence presentation and the
reviewer's decision path.

\fld{Potential} He works across product and engineering, so he can argue with
the constraint instead of designing around it. That combination is rare, and it
is what turns a technical result into a usable one.
\end{person}

%% ------------------------------------------------------ 2
\heading{2.}{Fraud and security}

\begin{person}{Abderrazak Hajji \ $\cdot$ \ Cybersecurity / Fraud Specialist}
\fld{Contact} +212 620 671 238 \ $\cdot$ \
\href{https://www.linkedin.com/in/abderrazak-hajji-45a2132bb/}{linkedin.com/in/abderrazak-hajji-45a2132bb}

\fld{Key capability} Cybersecurity, threat analysis, fraud detection.

\fld{Why him} My tests proved the system misses well-made forgeries. Deciding
which fraud patterns to prioritise, and where a false positive costs more than a
miss, requires someone who has seen real fraud.

\fld{Role} Define fraud scenarios, calibrate severity and the false-positive
threshold, challenge the detection rules adversarially.

\fld{Potential} He thinks like an attacker, which is the only reliable way to
test a detection tool. Someone whose instinct is to break the system is worth
more here than another builder.
\end{person}

%% ------------------------------------------------------ 3
\heading{3.}{Backend, cloud and MLOps}

\begin{person}{Adama \ $\cdot$ \ Backend / Cloud / MLOps Engineer}
\fld{Contact} +212 614 166 631 \ $\cdot$ \
\href{https://www.linkedin.com/in/startlingadama/}{linkedin.com/in/startlingadama}

\fld{Key capability} Backend engineering, DevOps, cloud infrastructure, AI systems.

\fld{Why him} The prototype has no rate limiting, no persistence and a payload
ceiling on its host. Each of these is a production problem, and each one blocks
a first paying customer.

\fld{Role} Build the API, the infrastructure and the deployment pipeline; make
the system scale and stay within cost.

\fld{Potential} He covers software, DevOps and AI at once --- a range that
usually takes three people. On a small team, that breadth is the difference
between shipping and stalling.
\end{person}

%% ------------------------------------------------------ clôture
\vspace{6pt}
\begin{note}{What this team would look like}
Four people covering detection, product, adversarial testing and production,
with me on the AI and detection side. Deliberately no second AI engineer --- the
prototype's failures were never a lack of modelling: they were a lack of
usability, domain knowledge and operational rigour.
\end{note}

\vspace{4pt}
{\footnotesize\color{muted}
Live prototype: \href{https://trust-doc-ai.vercel.app}{trust-doc-ai.vercel.app}
\hfill
Source: \href{https://github.com/Hajji-ahmed/TrustDoc-AI}{github.com/Hajji-ahmed/TrustDoc-AI}}

\end{document}
